\documentclass[10pt,a4paper]{amsart}
\usepackage[margin=19mm]{geometry}
\usepackage{amsmath,amssymb,mathtools}
\usepackage[scr=rsfs,cal=euler]{mathalfa}
\usepackage{xfrac}
\usepackage[unicode,hidelinks,bookmarks=false]{hyperref}
\newcommand{\ie}{i.e.\ }
\newcommand{\cf}{cf.\ }
\newcommand{\resp}{respectively\ }
\newcommand{\Subsection}{\subsection}
\providecommand{\nobreakdash}{\nobreak\hbox{-}}
\setlength{\emergencystretch}{2em}
\multlinegap0pt
\usepackage{amssymb}
\usepackage[normalem]{ulem}
\usepackage{tikz}
\newtheorem{lem}{Lemma}
\newtheorem{thm}{Theorem}
\newcommand\bK{{\mathbb{K}}}
\newcommand\bR{{\mathbb{R}}}
\newcommand\bZ{{\mathbb{Z}}}
\newcommand\inv{^{-1}}
\newcommand\pa{\partial}
\newcommand\rot{\operatorname{rot}}
\title{$\mathbb{K}$-framings and $\mathbb{K}$-quadratic forms on surfaces}
\author{Nariya Kawazumi}
\date{Source: arXiv:2604.27531v2 (CC BY 4.0)}
\begin{document}
\maketitle

\noindent\emph{Source and licence.} $\mathbb{K}$-framings and $\mathbb{K}$-quadratic forms on surfaces. Version arXiv:2604.27531v2 (CC BY 4.0), distributed by arXiv under the Creative Commons Attribution 4.0 International licence, http://creativecommons.org/licenses/by/4.0/, as recorded in the arXiv licence metadata for this exact version and shown on the abstract page https://arxiv.org/abs/2604.27531v2. Full licence notice: \texttt{licenses/arxiv-2604.27531-CC-BY-4.0.txt}.\par\smallskip

\noindent\textit{Editorial scope.} Counted unit: the article's thm thm:nontrivial (source0.tex lines 777--778). The article's complete original proof of it is delivered with it (source0.tex lines 779--823, one \texttt{proof} environment of 45 lines). No mathematics is written, repaired or re-derived: the statement and the proof are the article's own lines, in the article's own notation, and no retained line is substituted or re-flowed.  Retained uncounted as the source-local context the proof needs: lines 580--587; lines 758--764.  Nothing was omitted from any retained span, so the package carries no omission record.  No retained line contains a citation, so the wrapper carries no bibliography and no bibliography entry.  No retained line contains a figure environment, an image-inclusion command or drawing code, so no image is delivered and the package declares no asset dependency.  Editorial formatting only, in addition: an AMS \texttt{amsart} preamble that restates the article's own declarations and macros used by the retained lines, namely the environments \texttt{lem}, \texttt{thm} on separate counters, together with the standard packages those lines need.  The retained environments print in this document as Lemma 1, Theorem 1 in order of appearance, whereas the article prints the same items with the numbering of its own full document.  The complete native source of the article as distributed in the arXiv e-print for this exact version is bundled unchanged as \texttt{sources/199480/source0.tex}.
\begin{lem}[Poincar\'e-Hopf]\label{lem:PH}
We have 
\begin{equation}\label{eq:PH}
{\sum}^n_{j=0}\rot_\xi(\pa_j\Sigma) = \chi(\Sigma) \in \bK
\end{equation}
for any $\bK$-framing $\xi \in F(\Sigma; \bK)$. 
Here we endow each $\pa_j\Sigma$ with the orientation induced by that of the surface $\Sigma$.
\end{lem}

\begin{lem}\label{lem:DT}
Let $C \subset \Sigma$ be a simple closed curve. Then the value of $k_\xi$ at the right-handed Dehn twist $T_C$ along $C$ is given by
$$
k_\xi(T_C) = (\rot_\xi C)(\cdot [C]) \in H^1(\Sigma; \bK),
$$
where $\cdot [C]\in H^1(\Sigma, \pa \Sigma; \bZ)$ means the algebraic intersection with the homology class $[C] \in H_1(\Sigma; \bZ)$, i.e., the Poincar\'e duality of the class $[C]$. 
\end{lem}

\begin{thm}\label{thm:nontrivial} The class $k_\bK \in H^1(\Gamma_{g,n+1}; H^1(\Sigma_{g,n+1}; \bK))$ is nontrivial if and only if $g \geq 2$.
\end{thm} 

\begin{proof} For any simple closed curve $C \subset \Sigma$ and any $u \in H^1(\Sigma; \bK)$, the value $(\delta u)(T_C)$ of the coboundary $\delta u \in Z^1(\Gamma; H^1(\Sigma; \bK))$ at the Dehn twist $T_C$ is given by
\begin{equation}\label{eq:cob}
(\delta u)(T_C) = \langle u, [C]\rangle (\cdot [C]).
\end{equation}
In fact, for any loop $\gamma$ in $\Sigma$, 
we have $
\langle (\delta u)(T_C), [\gamma]\rangle =
\langle u\circ ({T_C}\inv) - u, [\gamma]\rangle =
\langle u, ({T_C}\inv)_*[\gamma] -[\gamma]\rangle=
\langle u, ([\gamma] \cdot [C])[C]\rangle=
\langle u, [C]\rangle ([\gamma] \cdot [C])
$.
\par
(I) The case $g=0$. If $n \leq 0$, it is trivial since $H^1(\Sigma_{0,n+1}; \bK) = 0$. So we suppose $n \geq 1$. Then the homology classes $\overrightarrow{\pa_j\Sigma}$, $1 \leq j \leq n$, and the fiber class $z$ generate
the homology group $H_1(U\Sigma; \bZ)$. The mapping class group $\Gamma = \Gamma_{0,n+1}$ acts trivially on these homology classes, and so on the cohomology group $H^1(U\Sigma; \bK)$. Hence, in the case $g =0$, we have 
\begin{equation}\label{eq:g0kxi}
k_\xi = 0: \Gamma \to H^1(\Sigma; \bK).
\end{equation}
\par
(II) The case $g=1$. By capping disks on all the boundary components, we can embed $\Sigma_{1,n+1}$ into the torus $\Sigma_{1,0}$. This induces a map $H^1(\Gamma_{1,0}; H^1(\Sigma_{1,0}; \bK)) \to H^1(\Gamma_{1,n+1}; H^1(\Sigma_{1,n+1}; \bK))$ which maps $k_{\bK}$ for $\Sigma_{1,0}$ to 
that for $\Sigma_{1,n+1}$, since any $\bK$-framing on $\Sigma_{1,0}$ induces a $\bK$-framing on $\Sigma_{1,n+1}$.  Here it should be remarked the torus is parallelizable. \par
Hence it suffices to show the vanishing of $k_{\bK}$ for $\Sigma_{1,0}$. 
Under the identification $\Sigma_{1,0} = \bR^2/\bZ^2$, we take two emmbeded loop $\alpha: t \in I \mapsto (t,0)\bmod{\bZ^2} \in \Sigma_{1,0}$ and $\beta: t \in I \mapsto (0, t)\bmod{\bZ^2} \in \Sigma_{1,0}$. %\par
Let $\xi \in F(\Sigma_{1,0}; \bK)$. 
From \eqref{eq:cob} we have $(\delta(\cdot[\gamma]))(T_C) = -([\gamma]\cdot[C])(\cdot [C])$ for any loop $\gamma$ in $\Sigma$. 
Hence we compute the coboundary $\delta u_\xi$ of the element $u_\xi := -  (\rot_\xi\beta)(\cdot [\alpha]) + (\rot_\xi\alpha)(\cdot [\beta]) \in H^1(\Sigma_{1,0}; \bK)$ as
$$
(\delta u_\xi)(T_C) = \left((\rot_\xi\beta)([\alpha]\cdot [C]) - (\rot_\xi\alpha)([\beta]\cdot [C]) \right)(\cdot [C]).
$$
In particular, we have $(\delta u_\xi)(T_\alpha) = (\rot_\xi\alpha)(\cdot [\alpha]) = k_\xi(T_\alpha)$ and $(\delta u_\xi)(T_\beta) = (\rot_\xi\beta)(\cdot [\beta]) = k_\xi(T_\beta)$ from Lemma \ref{lem:DT}. 
This implies $\delta u_\xi = k_\xi: \Gamma_{1,0} \to H^1(\Sigma_{1,0}; \bK)$,  
since the mapping class group $\Gamma = \Gamma_{1,0} \cong SL_2(\bZ)$ is generated by the Dehn twists $T_\alpha$ and $T_\beta$. This proves $k_\bK = 0$ in the case $g=1$. \par
(III) The case $g \geq 2$. Then we have a subsurface $S$ of $\Sigma$ which is diffeomorphic to a pair of pants, and each of whose boundary components is non-separating in the surface $\Sigma$. Let $C_k := \pa_kS$, $1 \leq k \leq 3$, be the three boundary components of the pair of pants $S$. 
Now assume that $k_\xi = \delta u$ for some $u \in H^1(\Sigma; \bK)$. 
From Lemma \ref{lem:DT} and \eqref{eq:cob}, we have $(\rot_\xi C_k)(\cdot [C_k])
= k_\xi(T_{C_k}) = (\delta u)(T_{C_k}) = \langle u, [C_k]\rangle (\cdot [C_k])$. 
Since $C_k$ is non-separating, we have $[\gamma_k]\cdot [C_k] = 1$ for some loop $\gamma_k$ in $\Sigma$. Hence we obtain $\rot_\xi C_k = \langle u, [C_k]\rangle$.
Consequently we obtain a contradiction 
$$
-1 = \chi(S) = {\sum}^3_{k=1}\rot_\xi C_k = \langle u, {\sum}^3_{k=1}[C_k]\rangle = 0
$$
from Lemma \ref{lem:PH} and $\sum^3_{k=1}[C_k] = 0 \in H_1(\Sigma; \bZ)$. 
This implies there is no $u \in H^1(\Sigma; \bK)$ satisfying $k_\xi = \delta u$, and 
proves $k_\bK \not= 0$ in the case $g\geq 2$. 
\end{proof}

\end{document}
